\section{Panoramica dell'Ontologia}
\label{sec:ontologia_overview}

\subsection{Dominio}

L'ontologia modella il dominio dei \textbf{videogiochi} nel periodo 2010--2026,
coprendo le seguenti entit\`a principali:
\begin{itemize}[nosep]
  \item \textbf{VideoGame} --- il prodotto videoludico con i suoi attributi (nome, data rilascio, punteggio Metacritic, paese d'origine);
  \item \textbf{Developer / Publisher} --- gli studi di sviluppo e le case editrici;
  \item \textbf{Genre} --- il genere (Action, RPG, Strategy, \dots{});
  \item \textbf{Platform} --- la piattaforma di esecuzione (PS5, PC, Switch, \dots{});
  \item \textbf{Character} --- i personaggi (giocanti e non);
  \item \textbf{Franchise} --- le serie e i sotto-franchise;
  \item \textbf{Award} --- i premi ricevuti;
  \item \textbf{GameEngine} --- il motore grafico utilizzato;
  \item \textbf{GameMode} --- la modalit\`a di gioco (single-player, multiplayer, co-op, MMO);
  \item \textbf{Review / GameEvent} --- recensioni ed eventi del ciclo di vita.
\end{itemize}

\subsection{Namespace e Separazione Schema/Istanze}

Il namespace dell'ontologia \`e:
\begin{center}
\texttt{http://www.videogame-ontology.org/ontology\#}
\end{center}
abbreviato in \texttt{vg:} in tutte le query SPARQL e nei riferimenti interni.

L'ontologia \`e strutturata in due file distinti che realizzano la classica
separazione schema/istanze:
\begin{itemize}[nosep]
  \item \texttt{videogames.owl} --- \textbf{schema}: contiene tutte le dichiarazioni
    di classi, propriet\`a, assiomi, disjointness e allineamenti esterni. Circa 934 triple;
  \item \texttt{videogames\_pruned\_2020.owl} --- \textbf{insieme di istanze di demo}: contiene le istanze
    (giochi, developer, publisher, generi, piattaforme, \dots{}) per il periodo
    2020--2026. Circa 745.000 triple, 68.700 giochi.
\end{itemize}

\subsection{Fonte Dati: Wikidata}

Tutti i dati provengono da \textbf{Wikidata} \citep{wikidata}, il knowledge graph
libero della Wikimedia Foundation. La popolazione avviene tramite query SPARQL
all'endpoint \url{https://query.wikidata.org/sparql}, mappando le propriet\`a Wikidata
rilevanti (identificatori P31, P136, P400, P577, P444, \dots{}) nelle corrispondenti
propriet\`a dell'ontologia.

\subsection{Pipeline di Popolamento}

La pipeline completa si articola in quattro fasi:

\begin{enumerate}[nosep]
  \item \texttt{populate\_wikidata.py} --- esegue 204 iterazioni (17 anni $\times$ 12 mesi)
    di query SPARQL verso Wikidata, raccogliendo giochi, metadati e relazioni.
    Tempo stimato: $\sim$5 ore con polite delay di 1 secondo tra le richieste.
  \item \textbf{OWL-RL closure via \texttt{owlrl}} --- applica la chiusura deduttiva
    OWL 2 RL per materializzare propriet\`a inverse, catene \texttt{subPropertyOf},
    e inferenze da \texttt{someValuesFrom} e \texttt{intersectionOf}. Da $\sim$1M a $\sim$1.4M triple.
  \item \textbf{Pruning} --- rimuove triple ridondanti o inutilizzate a runtime:
    \texttt{owl:sameAs} riflessivi, \texttt{gameDescription}, \texttt{officialWebsite}.
    Da $\sim$1.4M a $\sim$1.17M triple (completo) o $\sim$745k (demo 2020+).
  \item \texttt{enrich\_owl.py} --- gestisce gli assiomi DL-only (FunctionalProperty sulle
    datatype property e \texttt{hasKey}, cfr.\ Sezione~\ref{sec:reasoning}): li rimuove dalla
    copia di lavoro prima del reasoning e, con l'opzione \texttt{-{}-keep-dl-axioms}, li
    ripristina nel file serializzato.
\end{enumerate}

\subsection{Statistiche Riassuntive}

La Tabella~\ref{tab:ontostats} riassume le principali metriche dello schema dell'ontologia.

\begin{table}[h]
\centering
\caption{Statistiche dello schema dell'ontologia}
\label{tab:ontostats}
\begin{tabular}{lr}
\toprule
\textbf{Metrica}                & \textbf{Valore} \\
\midrule
Classi totali                   & 78              \\
Classi primitive                & 61              \\
Classi definite (equivalentClass) & 15            \\
Classi con vincolo di cardinalit\`a & 2           \\
Object Property                 & 31              \\
Datatype Property               & 26              \\
Property chain axioms           & 3               \\
Propriet\`a transitive          & 3               \\
Propriet\`a simmetriche         & 3               \\
Coppie inverse                  & 12              \\
Blocchi AllDisjointClasses      & 12              \\
Vincoli hasKey                  & 3               \\
Allineamenti equivalentClass    & 4 (Wikidata)    \\
Allineamenti equivalentProperty & 3 (Dublin Core) \\
Triple dello schema totali      & $\sim$ 934      \\
\bottomrule
\end{tabular}
\end{table}

Lo schema \`e espresso in formato RDF/XML; l'Appendice~\ref{sec:app_turtle} ne riporta
gli assiomi principali in Turtle. L'espressivit\`a DL e il profilo OWL 2 sono discussi
nella Sezione~\ref{sec:profilo_dl}, il reasoning nella Sezione~\ref{sec:reasoning}.
